\documentclass{article}
\PassOptionsToPackage{table,xcdraw}{xcolor}
\usepackage{xcolor}
\usepackage{bbm}
\usepackage{multirow}
\usepackage[english]{babel}

\usepackage[letterpaper,top=2cm,bottom=2cm,left=3cm,right=3cm,marginparwidth=1.75cm]{geometry}

\usepackage{amsmath}
\usepackage{graphicx}
\usepackage[colorlinks=true, allcolors=blue]{hyperref}
\usepackage{booktabs}
\usepackage{float}
\usepackage{caption}
\newcommand{\FrugalToolCall}{\textsc{FrugalToolCall}}

\newcommand{\eat}[1]{}

\usepackage{graphicx}
\usepackage{amssymb}
\usepackage{multirow}
\usepackage{bigstrut,morefloats}

\usepackage[resetlabels,labeled]{multibib}
\usepackage{caption}
\usepackage{subcaption}

\usepackage[most]{tcolorbox}

\newcommand{\E}{\mathbb{E}}

\usepackage{amsmath}
\usepackage{amsthm}
\usepackage{thmtools}

\theoremstyle{definition}

% algorithm2e conflicts with the algorithm+algpseudocode packages used by the
% profiler paper — replaced below.
%\usepackage[boxruled,linesnumbered]{algorithm2e}

\usepackage{algorithm}
\usepackage{algpseudocode}

\usepackage{array}

\usepackage[authoryear]{natbib}

% -------------------------------------------------------
% Profiler paper additions
% -------------------------------------------------------

\usepackage{xspace}
\usepackage[inline]{enumitem}

% Code listings (SQL and Python blocks in method.tex)
\usepackage{listings}
\lstdefinestyle{sqlstyle}{
  language=SQL,
  basicstyle=\footnotesize\ttfamily,
  keywordstyle=\color{blue!70!black}\bfseries,
  commentstyle=\color{gray},
  stringstyle=\color{red!60!black},
  breaklines=true,
  frame=single,
  numbers=none,
  xleftmargin=2pt,
  xrightmargin=2pt,
}
\lstdefinestyle{pystyle}{
  language=Python,
  basicstyle=\footnotesize\ttfamily,
  keywordstyle=\color{blue!70!black}\bfseries,
  commentstyle=\color{gray!70},
  stringstyle=\color{red!60!black},
  breaklines=true,
  frame=single,
  numbers=none,
  xleftmargin=2pt,
  xrightmargin=2pt,
}

% Row colors for attribution tiers table (xcolor table option covers \rowcolor)
\definecolor{highcol}{HTML}{1f4e79}
\definecolor{medcol}{HTML}{2e75b6}
\definecolor{ucol}{HTML}{bfbfbf}

% TODO / MEASURE markers — remove these before final submission
\usepackage[colorinlistoftodos]{todonotes}
\newcommand{\TODO}[1]{\todo[inline,color=yellow!40]{#1}}
\newcommand{\MEASURE}[1]{\todo[inline,color=orange!40]{MEASURE: #1}}

% Profiler convenience macros
\newcommand{\sys}{Operator Profiler\xspace}
\newcommand{\nsys}{\texttt{nsys}\xspace}
\newcommand{\ncu}{\texttt{ncu}\xspace}
\newcommand{\aten}[1]{\texttt{aten::#1}}
\newcommand{\figref}[1]{Figure~\ref{#1}}
\newcommand{\tabref}[1]{Table~\ref{#1}}
\newcommand{\secref}[1]{Section~\ref{#1}}
\newcommand{\algref}[1]{Algorithm~\ref{#1}}
\newcommand{\equref}[1]{Eq.~(\ref{#1})}
\newcommand{\etal}{\textit{et~al.}\xspace}
\newcommand{\eg}{\textit{e.g.,}\xspace}
\newcommand{\ie}{\textit{i.e.,}\xspace}

\usepackage{pifont}
\usepackage{microtype}
\usepackage{tikz}
\usepackage{pgfplots}
\pgfplotsset{compat=1.18}
\usetikzlibrary{arrows.meta,shapes,positioning,fit,backgrounds,calc,decorations.pathreplacing}

\title{Hardware-Attributed Operator Profiling for PyTorch}

\author{
Logan Chu$^{1}$ \quad
Dong Li$^{2}$ \\
\\
$^{1}$Yotta Labs \quad
$^{2}$Yotta Labs
}
\date{}

\renewcommand{\cite}[1]{\citep{#1}}

\begin{document}
\maketitle

\begin{abstract}
Framework profilers expose operator timing without hardware counters; GPU profilers expose
hardware counters without operator attribution. Bridging this gap manually is error-prone and
does not scale. We present \textbf{\sys}, a hardware attribution pipeline that automatically
links hardware metrics to PyTorch operators via three
complementary attribution paths: torch.profiler CUPTI correlation,
NVTX temporal enclosure with per-stream interval trees, and Inductor fusion-map
enrichment from debug artifacts. NVIDIA Nsight Compute (ncu) hardware counters are matched to NVIDIA Nsight Systems (nsys) kernel records
via invocation-order matching, avoiding timestamp joins across incompatible clock domains. A curated 20-counter metric set with
duration-weighted aggregation covers all hardware bottleneck axes, layer deduplication reduces \ncu replay time by a factor of $N/K$ for models with $N$ layers across $K$ unique structural classes, and
GPU clock locking controls the kernel-duration aggregates used for operator-level comparison.
On an NVIDIA RTX PRO 6000 Blackwell, \sys attributes 95--100\% of kernel
runtime for compiled workloads (GPT-2, SDPA Attention); black-box library backends
such as cuDNN RNN are correctly surfaced as ${>}$85\% unattributed rather than silently dropped.
Applied to profile-guided FX graph optimization, attributed profiles yield
\textbf{1.76$\times$--2.24$\times$} profiled-kernel-time speedups on the two
compiled optimization case studies; a third LSTM diagnostic case identifies
cuDNN re-dispatch as a structural fix rather than an FX graph rewrite.
\end{abstract}

\section{Introduction}
\label{sec:intro}

The throughput of a GPU workload is typically dominated by a small number of hardware bottleneck
classes: memory bandwidth saturation, Tensor Core under-utilization, register-pressure-limited
occupancy, and wave starvation from insufficient parallelism. The compute-bound vs.\ memory-bound
taxonomy that Williams \etal~\cite{williams2009roofline} formalized for multicore architectures
remains difficult to apply at the operator level in PyTorch, not because the hardware data is
absent, but because no tooling bridges it to the framework abstraction.

\paragraph{The attribution gap.}
Framework profilers (\texttt{torch.profiler}, \nsys)
expose operator timing without hardware counters; \ncu exposes hardware
counters without operator attribution. Bridging the two requires matching records across tools
that run in separate passes, use incompatible clock domains (CUPTI GPU timestamps vs.\ \ncu
injection-mode timestamps), and produce outputs with no shared stable key. No existing tool was designed to bridge all three (\secref{sec:related}).

\paragraph{PyTorch compilation stack.}
In PyTorch's default eager execution, each \texttt{aten::} operator dispatches individually
through the operator dispatcher at runtime. Under TorchInductor, operators
for which a hand-tuned vendor library implementation exists---such as cuBLAS for matrix
multiplication or xFormers for attention---are emitted as extern calls that
preserve this per-operator dispatch. Operators without such external counterparts---elementwise
operations, reductions, normalization, and layout transforms---are fused in groups into
Triton~\cite{triton2019} kernels, JIT-compiled to CUDA, where each kernel may represent
multiple operators and individual \texttt{aten::} dispatch does not occur at runtime.

\paragraph{Our approach.}
We observe that reliable attribution does not require timestamp matching.  Because \ncu replays
kernels in the same execution order they appear in the \nsys trace, a kernel's
\texttt{(name, invocation\_index)} pair is a stable, clock-immune identifier.  \sys combines
three complementary attribution paths: NVTX temporal enclosure for \texttt{aten::}-dispatched
kernels, invocation-order matching for \ncu counter assignment, and Inductor debug-artifact
parsing for Triton-fused kernels with no NVTX bracket.  \figref{fig:overview} shows the
complete end-to-end pipeline.

\begin{figure}[H]
\centering
\begin{tikzpicture}[
  font=\small,
  box/.style={rounded corners=3pt, draw, align=center, minimum height=1.8em, inner sep=5pt},
  toolbox/.style={box, fill=blue!12, draw=blue!50!black},
  attrbox/.style={box, fill=orange!15, draw=orange!70!black},
  outbox/.style={box, fill=green!12, draw=green!50!black},
  pgobox/.style={box, fill=red!10, draw=red!60!black},
  arr/.style={-{Stealth[length=5pt]}, thick},
  darr/.style={-{Stealth[length=5pt]}, thick, dashed, gray},
  lbl/.style={font=\scriptsize\itshape, text=gray!70!black},
]

%% ---- Row 1: Profiling stage ----
\node[toolbox] (workload) at (0,0)   {\texttt{workload.py}};
\node[toolbox] (nsys)     at (2.7,0) {\nsys capture};
\node[toolbox] (sqlite)   at (5.5,0) {SQLite export};
\node[toolbox] (ncu)      at (8.2,0) {\ncu replay};

\draw[arr] (workload) -- (nsys);
\draw[arr] (nsys)     -- (sqlite);
\draw[arr] (sqlite)   -- (ncu);

%% ---- Row 2: Attribution engine ----
\node[attrbox] (iom)    at (8.2,-1.65) {Invocation-Order\\[-1pt]Matching};
\node[attrbox] (sit)    at (5,-1.65) {StreamIntervalTree\\[-1pt]NVTX enclosure};
\node[outbox]  (prof)   at (1.5,-1.65) {\texttt{profile.json}\\[-1pt](20 counters)};

\draw[arr] (sqlite) -- (iom);
\draw[arr] (ncu)    -- (iom);
\draw[arr] (iom)    -- (sit);
\draw[arr] (sit)    -- (prof);


%% ---- Row 3: PGO ----
\node[pgobox] (diag)    at (2.7,-3.3)  {Diagnosis\\[-1pt]Engine};
\node[pgobox] (passes)  at (5.5,-3.3)  {Optimization\\[-1pt]Patterns};
\node[pgobox] (compile) at (8.2,-3.3)  {\texttt{torch.compile}\\[-1pt]+ Inductor};
\node[outbox] (opt)     at (10.8,-3.3) {Optimized\\[-1pt]model};

\draw[arr] (prof)    -- (diag);
\draw[arr] (diag)    -- (passes);
\draw[arr] (passes)  -- (compile);
\draw[arr] (compile) -- (opt);


%% ---- Feedback loop ----
\draw[darr] (opt.south) -- ++(0,-0.45) -| (workload.south)
  node[midway, below, lbl] {closed-loop validation };

\end{tikzpicture}
\caption{\textbf{\sys system overview.}  Two profiling stages combine to produce \texttt{profile.json}: \nsys
  captures the kernel timeline and \ncu captures hardware counters. The attribution engine bridges the two tools via invocation-order matching and per-stream NVTX interval trees, constituting the core contribution of \sys.  The lowest ``Optimization'' row shows a downstream application of the attributed profile: hardware-evidence-motivated
  FX graph passes validated via re-profiling; it is not part of the attribution pipeline itself.}
\label{fig:overview}
\end{figure}

\paragraph{Contribution.}
We present \textbf{\sys}, a hardware attribution pipeline that automatically links \ncu hardware
counters to PyTorch \texttt{aten::} operators via three complementary paths: NVTX temporal
enclosure, invocation-order matching, and Inductor fusion-map enrichment. Clock-domain joins are
avoided entirely; unattributed kernels are surfaced explicitly, never dropped. The pipeline
collects a curated 20-counter metric set spanning all roofline bottleneck axes via
application-mode \ncu replay, with architecture-specific fallbacks for Ampere, Hopper, and
Blackwell, and supports layer deduplication to reduce replay time by
$O(\text{num\_duplicate\_layers})$ for models with repeated structure.

\section{Related Work}
\label{sec:related}

The tools closest to \sys fall short along three primary dimensions: operator-level timing
without counters, kernel-level counter collection without operator identity, and
compilation-level optimization without counter-verified feedback. Adjacent areas---fused kernel
libraries and profile-guided optimization---are surveyed below for completeness.

\paragraph{Operator-level profilers.}
\texttt{torch.profiler} and its Kineto backend \cite{kineto2021} record wall-clock durations and
Chrome traces at the \texttt{aten::} operator granularity~\cite{pytorch2019} and expose CUDA activity correlation
via CUPTI External IDs---giving kernel-level \emph{timing} attribution but no access to hardware
performance counters. Practitioners can identify \emph{which} operator is slow, but cannot
determine \emph{which hardware resource} it is contending for---the distinction that determines
what optimization is warranted. Without hardware
counter data at the operator level, optimization choices remain heuristic.
NVIDIA DLProf~\cite{dlprof} is the closest NVIDIA-specific predecessor: it uses NVTX markers
to correlate CPU and GPU time with model operations, reports at operation, layer, and kernel
granularity, and includes Tensor Core usage diagnostics. DLProf remains timing-centric:
its GPU-side data comes from \nsys, and Tensor Core utilization is inferred from kernel names.
\sys adds the missing counter-attribution layer. It joins application-mode \ncu hardware
counter rows to \nsys kernel records through invocation-order matching, then assigns those
counter-enriched kernel records to PyTorch operators, including Inductor-era fused kernels.
This produces per-operator hardware metrics such as measured Tensor Core activity,
occupancy, memory throughput, register pressure, and cache behavior rather than only
operation timing or kernel-name-derived Tensor Core eligibility.

\paragraph{Kernel-level profilers.}
NVIDIA Nsight Systems (\nsys) \cite{nsightsystems} captures the full CUDA kernel timeline with
NVTX annotations at high throughput, but produces no operator-level view: every kernel appears by
its internal CUDA or Triton name (\eg \texttt{cutlass\_80\_simt\_sgemm\_128x32\_tn}), with no
indication of which PyTorch operator dispatched it. NVIDIA Nsight Compute (\ncu)
\cite{nsightcompute} collects per-kernel hardware counters at full precision---occupancy, DRAM
bytes, Tensor Core utilization, cache hit rates---but its output is keyed on kernel names and
invocation indices in the replay timeline, again with no awareness of the operator graph above.
NVIDIA DCGM and Dynolog \cite{dynolog2023} operate at cluster monitoring granularity and do not
perform per-kernel or per-operator attribution. Taken together, these tools collectively address
hardware counter collection and kernel-timeline capture but leave kernel-to-operator attribution entirely unresolved.

\paragraph{Deep learning compilers and optimizers.}
TorchDynamo and TorchInductor \cite{pytorch2024} compile PyTorch programs via FX graph capture~\cite{torchfx2022} to Triton kernels
using heuristic fusion rules (pointwise chaining, matmul epilogue fusion), but fusion decisions
are made without hardware counter feedback: whether a given fusion improved Tensor Core
utilization is not a question the standard Inductor toolchain can answer programmatically. TVM \cite{tvm2018} and Halide \cite{halide2013} optimize tensor programs via
measurement-based auto-tuning and schedule/algorithm separation respectively; both optimize
\emph{latency} directly rather than using hardware counters to identify the bottleneck class
driving that latency. ONNX Runtime\footnote{\url{https://onnxruntime.ai}} and NVIDIA
TensorRT\footnote{\url{https://developer.nvidia.com/tensorrt}} perform layer fusion and
quantization behind black-box interfaces that cannot be extended for custom architectures or
augmented with attribution data.

\paragraph{Fused kernel libraries.}
FlashAttention \cite{flashattention2022,flashattention2024} and xFormers \cite{xformers2022}
provide hand-optimized fused attention kernels that eliminate DRAM round-trips by keeping
intermediate tensors in on-chip SRAM. These represent the class of optimization that
attributed hardware counter evidence should motivate rather than apply unconditionally; what
is absent is a profiling layer that identifies, from measured per-operator counters, which
variant is warranted for a given workload.

\paragraph{Profile-guided optimization.}
Classic PGO in C/C++ compilers (GCC, LLVM) collects branch frequency profiles to guide inlining
and layout decisions \cite{autofdo2016}. Applying the same principle at the GPU operator level requires hardware counter data attributed
to specific operators. To our knowledge, \sys is the first published system to combine
\nsys kernel traces, application-mode \ncu hardware counters, and PyTorch operator
attribution through invocation-order matching, and to use the resulting attributed profile
to drive hardware-informed FX graph optimization.

\section{Attribution Pipeline}
\label{sec:attribution}

\sys ingests two data sources produced for the same workload run: (1) an \nsys SQLite export
containing CUDA kernel activity and NVTX range annotations, and (2) \ncu CSV output containing
per-kernel hardware counter measurements. The pipeline produces a hardware-attributed
\texttt{OperatorProfile}: a JSON document mapping each \texttt{aten::} operator to its constituent
kernels, their hardware metrics, and the confidence tier of each attribution.

\subsection{Data Ingestion}
\label{sec:ingestion}

\textbf{Nsight Systems export.}  Running the workload under
\texttt{nsys profile --trace=cuda,nvtx} produces a \texttt{.nsys-rep} file, exported to SQLite
via \texttt{nsys export --type=sqlite}. \sys queries three tables:

\begin{lstlisting}[style=sqlstyle,caption={Kernel and NVTX range extraction from the \nsys SQLite export.}]
-- CUDA kernel activity (GPU-domain timestamps)
SELECT k.correlationId, s.value AS kernel_name,
       k.start AS start_ns, k.end AS end_ns,
       k.streamId, k.deviceId,
       k.gridX, k.gridY, k.gridZ,
       k.blockX, k.blockY, k.blockZ,
       r.globalTid AS host_tid,
       r.start      AS cpu_launch_start_ns
FROM   CUPTI_ACTIVITY_KIND_KERNEL k
JOIN   StringIds s ON s.id = k.shortName
LEFT JOIN CUPTI_ACTIVITY_KIND_RUNTIME r
       ON r.correlationId = k.correlationId
ORDER  BY k.start ASC;

-- NVTX range annotations (CPU-domain timestamps)
SELECT n.text, n.start AS start_ns, n.end AS end_ns,
       n.nestingLevel, n.globalTid AS stream_id
FROM   NVTX_EVENTS n
WHERE  n.end IS NOT NULL
ORDER  BY n.start ASC;
\end{lstlisting}

The \texttt{cpu\_launch\_start\_ns} field from the RUNTIME table is the host-side timestamp when
the CUDA kernel launch API returned, providing a CPU-domain anchor for NVTX bracket queries
without conflating it with the GPU-domain kernel execution timestamp (\secref{sec:clockdomain}).

\textbf{Nsight Compute export.}  For the full set of kernels in the \nsys trace, \sys invokes
a single application-mode replay pass per counter group (4--8 passes total), collecting all
20 metrics simultaneously:
\begin{lstlisting}[style=pystyle]
ncu --replay-mode application \
    --metrics <comma_separated_list> \
    --export <output.ncu-rep> \
    [--kernel-name <filter>] \
    <replay_script.py>
\end{lstlisting}
Application-mode replays the workload once per counter group regardless of the number of
unique kernel names, making profiling cost $O(\text{counter groups})$ rather than
$O(\text{unique kernels})$---a 10--50$\times$ efficiency gain for multi-layer models.
The optional \texttt{--kernel-name} flag filters to a specific kernel for targeted
re-profiling.  The resulting \texttt{.ncu-rep} is exported to CSV
(\texttt{ncu --import --csv}), producing one row per kernel invocation per metric.

\subsection{GPU Clock Locking for Reproducible Duration Comparisons}
\label{sec:clocklocking}

Per-operator durations in \texttt{profile.json} are derived from \nsys CUPTI GPU-domain
timestamps.  These durations are subject to a subtle confound: a GPU's boost clock floats
with die temperature and power draw, typically varying 5--15\% across captures separated
by minutes or hours.  An optimized workload dissipates less power than the baseline (fewer
FLOPs on the SIMT path), so its capture may run at a \emph{higher} sustained clock,
inflating the apparent speedup.  Conversely, thermal throttling during a long baseline
capture deflates the baseline.  Both effects corrupt the $\text{baseline}/\text{optimized}$
duration ratio even when no hardware change was made.

\textbf{Probe-and-lock.}
\sys pins the SM and memory clocks for the duration of the \nsys capture.  The default
target is \emph{probe-and-lock}: a synthetic matrix-multiply load runs for ${\sim}1.5\,\text{s}$
(with a $1\,\text{s}$ warmup to reach thermal steady state), and the median clock observed
during the probe is selected as the lock target.  Locking to the \emph{sustained} frequency
(rather than the peak boost clock, which the GPU cannot maintain under load) guarantees the
lock is genuinely held throughout the capture.

\textbf{Cache and reuse.}
The probed clock is written to \texttt{profiler\_output/.gpu\_clock\_lock.json} with a
6-hour TTL per GPU index.  The baseline capture probes once; the optimized capture reuses
the cached clock.  Both captures therefore execute at exactly the same frequency, making
their duration ratio clock-immune by construction---the same guarantee that invocation-order
matching provides for the hardware counter join.

\textbf{Implementation.}
Clocks are set via \texttt{nvidia-smi -lgc/-lmc} and released through a context manager
that registers a \texttt{SIGTERM} handler, guaranteeing a reset even on exception, timeout,
or Ctrl-C:
\begin{lstlisting}[style=pystyle]
with gpu_clocks_locked(target="probe", gpu_index=0,
                       use_sudo=True, cache_dir=cache_dir):
    # nsys profile ... run_workload.py runs here
    ...
\end{lstlisting}
Clock setting uses \texttt{sudo -n} (non-interactive); if permission is unavailable, \sys
logs a \texttt{WARNING} and proceeds at dynamic clocks rather than aborting.

\textbf{Scope.}
Clock locking applies to the \nsys capture phase only.  The \ncu replay phase is
unaffected: \ncu already self-locks to its own base clock before each replay pass, and
the hardware counter values it reports are independent of the GPU's runtime SM frequency.

\subsection{The Clock Domain Problem and Invocation-Order Matching}
\label{sec:clockdomain}

The fundamental obstacle in cross-tool profiling is that \nsys and \ncu observe the same kernel
execution through incompatible timing mechanisms. \nsys records CUPTI GPU-domain timestamps from
the hardware's free-running clock. \ncu replays kernels in an isolated profiling pass with injected
performance counter reads; its internal invocation clock reflects replay-mode execution, not the
original capture clock. On a single GPU, these two clocks routinely diverge by 10--100~$\mu$s per
kernel due to replay serialization and PMU injection overhead. A naive timestamp join would
incorrectly associate hardware counter measurements with the wrong operator---especially dangerous
for operators with similar execution times.

\textbf{Invocation-Order Matching (Algorithm~\ref{alg:iom}).}  The key insight is that \ncu
replays kernels in the same execution order they appear in the \nsys timeline: the $i$-th
invocation of kernel $K$ in the \nsys trace corresponds to the $i$-th row for kernel $K$ in the
\ncu CSV output.  We exploit this ordering to assign hardware metrics without any timestamp
comparison:

\begin{algorithm}[t]
\caption{Invocation-Order Matching (\nsys $\times$ \ncu)}
\label{alg:iom}
\begin{algorithmic}[1]
\Require \texttt{nsys\_kernels}: list of \texttt{KernelRow} sorted by \texttt{start\_ns} per stream
\Require \texttt{ncu\_rows}: list of \texttt{NcuRow} sorted by invocation index per kernel name
\Ensure  Each \texttt{KernelRow} assigned a \texttt{KernelMetrics} object
\State $\textit{counter} \gets \text{defaultdict}(\text{int})$
  \Comment{per kernel name}
\For{$k \in \texttt{nsys\_kernels}$}
  \State $i \gets \textit{counter}[k.\textit{name}]$
  \If{$i \ge |\texttt{ncu\_rows}[k.\textit{name}]|$}
    \State \textbf{raise} \texttt{InvocationCountMismatchError}($k.\textit{name}$, $i$)
    \Comment{iteration count mismatch between captures}
  \EndIf
  \State $k.\textit{metrics} \gets \texttt{ncu\_rows}[k.\textit{name}][i]$
  \State $\textit{counter}[k.\textit{name}] \mathrel{+}= 1$
\EndFor
\end{algorithmic}
\end{algorithm}

This matching is correct when the workload is \emph{deterministic}: the same kernel names must
appear in the same counts and the same execution order in both the \nsys capture and the \ncu
replay.  Concretely, models with conditional control flow, dynamic dispatch, or
\texttt{cudaMallocAsync} workspace kernels can violate this invariant without producing a shape
error.  \sys validates that per-kernel invocation counts match between the two runs before
issuing counter assignments; a count mismatch on any kernel name raises
\texttt{InvocationCountMismatchError} and aborts attribution for that workload.  Shape validation
before the replay step enforces the shape portion of this invariant
(\secref{sec:edgecases}).  \figref{fig:iom} illustrates why timestamp joins fail and how
invocation-order matching avoids the problem.

\begin{figure}[t]
\centering
\begin{tikzpicture}[
  font=\small,
  arr/.style={-{Stealth[length=4pt]}, thick},
  karr/.style={-{Stealth[length=4pt]}, thick, red!70!black},
  garr/.style={-{Stealth[length=4pt]}, thick, green!50!black},
  kern/.style={draw, rounded corners=2pt, fill=blue!10, draw=blue!60!black,
               minimum width=1.6cm, minimum height=1.4em, align=center, font=\footnotesize},
  kernncu/.style={draw, rounded corners=2pt, fill=purple!10, draw=purple!60!black,
               minimum width=1.6cm, minimum height=1.4em, align=center, font=\footnotesize},
  tl/.style={-{Stealth[length=5pt]}, thick, gray},
  lbl/.style={font=\scriptsize, text=gray!80!black},
  hl/.style={font=\footnotesize\bfseries},
  darr/.style={-{Stealth[length=4pt]}, thick, dashed, gray!50},
]

%% ---- (a) Naive timestamp join ----
\node[hl] at (4.5, 0.7) {(a) Naive timestamp join --- \textcolor{red!70!black}{incorrect}};

\draw[tl] (0.0, 0) -- (9.8, 0) node[right,lbl] {nsys (GPU clock)};
\node[lbl, left=2pt] at (0.0, 0) {\small nsys};

\node[kern] (k1n) at (1.1, 0)  {$K_1$\;\tiny t=10};
\node[kern] (k2n) at (4.0, 0)  {$K_2$\;\tiny t=45};
\node[kern] (k3n) at (7.1, 0)  {$K_3$\;\tiny t=80};

\draw[tl] (0.0,-1.25) -- (9.8,-1.25) node[right,lbl] {ncu (replay clock)};
\node[lbl, left=2pt] at (0.0,-1.25) {\small ncu};

\node[kernncu] (k1c) at (2.3,-1.25)  {$K_1$\;\tiny t=35};
\node[kernncu] (k2c) at (5.1,-1.25)  {$K_2$\;\tiny t=62};
\node[kernncu] (k3c) at (8.1,-1.25)  {$K_3$\;\tiny t=95};

\draw[darr] (k1c.north) -- (k1n.south);
\draw[karr] (k2c.north) -- (k1n.south)
  node[near start, left=3pt, lbl, text=red!70!black] {\textbf{\ding{55}} wrong match};
\draw[darr] (k3c.north) -- (k3n.south);

\draw[decorate,decoration={brace,amplitude=5pt},gray!60]
  (k1n.south east)++(0.1,0) -- (k1c.north east)++(0.1,0)
  node[midway, right=6pt, lbl, align=left] {clock\\drift};

%% ---- (b) IOM ----
\node[hl] at (4.5,-2.15) {(b) Invocation-order matching (IOM) --- \textcolor{green!50!black}{correct}};

\draw[tl] (0.0,-2.9) -- (9.8,-2.9) node[right,lbl] {nsys order};
\node[lbl, left=2pt] at (0.0,-2.9) {\small nsys};
\node[kern] (k1ni) at (1.4,-2.9) {$K_1$\;\tiny $i{=}0$};
\node[kern] (k2ni) at (4.0,-2.9) {$K_2$\;\tiny $i{=}1$};
\node[kern] (k3ni) at (6.7,-2.9) {$K_3$\;\tiny $i{=}2$};

\draw[tl] (0.0,-4.1) -- (9.8,-4.1) node[right,lbl] {ncu order};
\node[lbl, left=2pt] at (0.0,-4.1) {\small ncu};
\node[kernncu] (k1ci) at (1.4,-4.1) {$K_1$\;\tiny $i{=}0$};
\node[kernncu] (k2ci) at (4.0,-4.1) {$K_2$\;\tiny $i{=}1$};
\node[kernncu] (k3ci) at (6.7,-4.1) {$K_3$\;\tiny $i{=}2$};

\draw[garr] (k1ci.north) -- (k1ni.south) node[midway,right=1pt,lbl,text=green!50!black] {\textbf{\ding{51}}};
\draw[garr] (k2ci.north) -- (k2ni.south) node[midway,right=1pt,lbl,text=green!50!black] {\textbf{\ding{51}}};
\draw[garr] (k3ci.north) -- (k3ni.south) node[midway,right=1pt,lbl,text=green!50!black] {\textbf{\ding{51}}};

\node[lbl, align=center] at (4.9,-4.75)
  {Match key: \texttt{(kernel\_name, invocation\_index)} --- no timestamp comparison required};

\end{tikzpicture}
\caption{\textbf{Why timestamp joins fail and how IOM avoids them.}
  (a)~The nsys GPU-domain clock and the \ncu replay clock diverge by $>$10\,$\mu$s; a
  nearest-timestamp join misattributes $K_2$'s hardware counters to $K_1$.
  (b)~IOM matches kernels by ordinal position within each \texttt{(kernel\_name,
  invocation\_index)} sequence, making attribution clock-immune by construction.}
\label{fig:iom}
\end{figure}

\subsection{StreamIntervalTree Attribution}
\label{sec:intervaltree}

For each unique \texttt{(stream\_id, device\_id)} pair observed in the \nsys trace, \sys
constructs a \texttt{StreamIntervalTree}: a sorted list of NVTX ranges indexed by
\texttt{start\_ns}, implemented using Python's \texttt{sortedcontainers.SortedList} for
$O(\log n)$ insertion and $O(k \log n)$ enclosure queries ($k$ = number of enclosing ranges).

\begin{algorithm}[t]
\caption{NVTX Enclosure Attribution (MEDIUM Confidence)}
\label{alg:nvtx}
\begin{algorithmic}[1]
\Require Kernel row $k$ with \texttt{cpu\_launch\_start\_ns}, \texttt{stream\_id}
\Require \texttt{forest}: map from \texttt{(stream\_id, device\_id)} to \texttt{SortedList}
\Ensure Attribution at MEDIUM confidence, or fall through to heuristic
\State $T \gets k.\textit{cpu\_launch\_start\_ns}$
\State $\textit{ranges} \gets$ \texttt{forest}[$k.\textit{stream\_id}$].query\_enclosing($T$)
  \Comment{all NVTX ranges $[s,e]$ with $s \le T \le e$}
\If{$|\textit{ranges}| = 0$}
  \State \textbf{fall through} to Inductor fusion enrichment (\secref{sec:heuristic}) or UNATTRIBUTED
\ElsIf{$|\textit{ranges}| = 1$}
  \State $k.\textit{operator} \gets \textit{ranges}[0].\textit{text}$,
         \textit{confidence} $\gets$ MEDIUM
\Else
  \State $k.\textit{operator} \gets \arg\max_r r.\textit{nesting\_level}$
         \Comment{innermost enclosing range; ties broken by latest \texttt{start\_ns}}
  \State $k.\textit{is\_fused} \gets \top$;
         $k.\textit{all\_enclosing} \gets \textit{ranges}$
         \Comment{fused kernel spans multiple operators}
\EndIf
\end{algorithmic}
\end{algorithm}

The \texttt{cpu\_launch\_start\_ns} field is used as the query point rather than the GPU
\texttt{start\_ns}, because NVTX ranges are CPU-side events and the async gap between host launch
and GPU execution makes GPU timestamps unreliable for enclosure queries.

\subsection{Inductor Fusion Enrichment}
\label{sec:heuristic}

When NVTX enclosure fails (no enclosing range exists at the query point), \sys applies a
post-attribution enrichment pass that recovers ground-truth fused operator identity from Inductor
debug artifacts, rather than guessing from kernel names.

\textbf{Inductor debug extraction.}  When \texttt{torch.compile(backend="inductor")} runs with
debug mode enabled, Inductor writes \texttt{output\_code.py} artifacts containing the generated
Triton kernel source.  Each kernel call site is preceded by a structured comment:

\begin{lstlisting}[style=pystyle]
# Topologically Sorted Source Nodes: [...],
# Original ATen: [aten.relu, aten.addmm]
def triton_poi_fused_relu_addmm_0(in_ptr0, ...):
    ...
\end{lstlisting}

\sys's \texttt{InductorFusionExtractor} parses these comments via the regex
\texttt{\# .+Original~ATen:\textbackslash{}s*\textbackslash{}[([\textasciicircum{}\textbackslash{}]]+)\textbackslash{}]},
producing a mapping $\{$\textit{kernel\_short\_name} $\to$ $[$\texttt{aten::op}_1,
\texttt{aten::op}_2, $\ldots]$$\}$ that contains the exact set of aten operators fused into
each kernel.

\textbf{Enrichment as a post-attribution pass.}  This map is applied \emph{after} the NVTX
enclosure pass and never overrides a valid TORCH\_PROFILER or NVTX attribution:
\begin{enumerate}[leftmargin=1.5em]
  \item For kernels that remain UNATTRIBUTED after NVTX enclosure, a lookup in the fusion map
    upgrades them to \texttt{INDUCTOR\_FUSION} attribution at MEDIUM confidence, with
    \texttt{source\_operators} populated from the \texttt{Original~ATen} comment.
  \item For already-attributed kernels (TORCH\_PROFILER or NVTX), the enrichment augments the
    \texttt{fused\_ops} field and sets \texttt{is\_fused~=~True} if multiple operators are present.
\end{enumerate}

Kernels matching neither NVTX enclosure nor the Inductor fusion map remain in
\texttt{unattributed\_kernels[]} with \texttt{attribution\_method = "unattributed"}.  For eager-mode
workloads or cuDNN-backed operators (\eg \texttt{nn.LSTM} via \texttt{aten::\_cudnn\_rnn}), which
produce neither NVTX brackets nor Inductor debug output, unattributed rates of 20--85\% are expected
and are surfaced explicitly in the profile output.

\subsection{Confidence Tiers}
\label{sec:tiers}

\begin{table}[t]
\centering
\caption{Attribution methods in \sys.  Coverage percentages are measured over kernel
  \emph{runtime} (not kernel count); the UNATTRIBUTED rate varies widely with compile mode and
  operator backend.  Ranges shown are typical across PyTorch-native workloads; see
  \tabref{tab:coverage} for per-workload values.}
\label{tab:tiers}
\small
\setlength{\tabcolsep}{4pt}
\begin{tabular}{p{2.0cm}p{2.6cm}p{4.8cm}p{1.6cm}}
\toprule
\textbf{Method} & \textbf{Evidence} & \textbf{Trigger Condition} & \textbf{Coverage} \\
\midrule
\rowcolor{highcol!15}
HIGH & torch.profiler IOM & Kineto trace available; kernel count agrees with dispatch count & 35--45\% \\
\rowcolor{medcol!15}
MEDIUM (NVTX) & NVTX enclosure (interval tree) & Kernel launch falls within an \texttt{aten::} NVTX range & 45--55\% \\
\rowcolor{medcol!10}
MEDIUM (Inductor) & Inductor \texttt{output\_code.py} comments & Kernel name in Inductor fusion map; NVTX enclosure absent & 5--15\% \\
\rowcolor{ucol!15}
UNATTRIBUTED & None & cuDNN-internal kernels (LSTM, cuDNN conv.); eager Triton without NVTX & 1--85\% \\
\bottomrule
\end{tabular}
\end{table}

\subsection{Edge Cases}
\label{sec:edgecases}

\sys explicitly handles the following attribution edge cases:

\begin{enumerate}[leftmargin=1.5em]

\item \textbf{Clock domain mismatch.} Invocation-order matching (\algref{alg:iom}) makes
  attribution clock-immune by construction.

\item \textbf{CUDA Graph replay.} During graph capture, operators emit NVTX ranges and kernels are
  attributed normally.  During replay, new GPU timestamps with no NVTX correspondence are generated.
  \sys detects replay by the absence of NVTX events in the replay window and replays the attribution
  cache from the capture phase.

\item \textbf{Multi-stream execution.} Each \texttt{(stream\_id, device\_id)} pair has its own
  \texttt{StreamIntervalTree}, preventing cross-stream false matches.

\item \textbf{JIT warm-up kernels.} Kernels launched before the first NVTX range (JIT compilation
  artifacts) are flagged using a duration-outlier heuristic ($> 10\times$ median kernel duration)
  and excluded from operator mapping.

\item \textbf{Async kernel launch.} NVTX push/pop occurs on the host before the kernel executes on
  the GPU.  The RUNTIME table's \texttt{cpu\_launch\_start\_ns} field provides a host-side anchor
  that lies within the NVTX range, whereas the GPU \texttt{start\_ns} may not.

\item \textbf{Dynamic shapes.} \sys validates that input shapes at \ncu replay match those recorded
  during \nsys capture; a \texttt{ShapeMismatchError} is raised otherwise.

\item \textbf{Fused kernels.} When multiple NVTX ranges enclose a single kernel (Triton-fused
  cross-operator), \sys sets \texttt{is\_fused=True} and stores all enclosing ranges in
  \texttt{all\_enclosing\_ranges} for downstream analysis.

\item \textbf{ncu timing independence.} \ncu timestamps are internal to each replay run and are
  discarded entirely; only metric values are extracted from \ncu output.

\end{enumerate}

\subsection{Layer Deduplication}
\label{sec:dedup}

For models with repeated structure (\eg transformer decoder blocks), profiling and
compiling every layer independently costs $O(N)$ ncu replay time.
\texttt{UniqueSubgraphRegistry} detects structurally identical FX subgraphs by
computing a canonical hash of each partition's node types and connectivity.  Only
one representative per equivalence class is compiled and profiled; metrics and
compiled callables are propagated to all structural duplicates by positional index
within each partition group.

The deduplication is transparent to the outer profiling loop: each duplicate
partition's forward pass is wrapped with an NVTX range tagged
\texttt{layer::duplicate::<label>}, and the partition equivalence map (written to
\texttt{<prefix>.part.json} during the correlation pass) is passed to
\texttt{KernelProfileConfig} so the \ncu orchestrator skips duplicate-partition
kernels during replay and propagates metrics from the unique representative.

\textbf{Impact.}  For GPT-2 (12 identical transformer blocks),
deduplication reduces \ncu replay time from ${\sim}$30--45\,min to
${\sim}$3--5\,min---a 12$\times$ improvement.  Layer deduplication is activated
automatically by \texttt{run\_workload.py} when structurally repeated partitions
are detected.

\textbf{Warmup-gating fix.}  An important correctness invariant: the
\texttt{layer::unique::} and \texttt{layer::duplicate::} NVTX ranges must only be
emitted during the \emph{measured} capture window, not during warmup iterations.
\texttt{run\_workload.py} gates all layer-partition NVTX emission behind an
\texttt{\_NVTX\_STATE[`active']} flag that is set only after warmup completes.
Without this guard, warmup-phase kernels appear in the nsys trace inside
layer-partition ranges, causing \texttt{ManifestBuilder} to misattribute them to
operators rather than classifying them as pre-NVTX initialization---inflating
attributed kernel counts and biasing operator durations.

% -------------------------------------------------------
\section{Metric Set and Aggregation}
\label{sec:metrics}

\subsection{Counter Selection Rationale}

Collecting all 90+ \ncu counters is expensive (multiple replay passes) and produces profiles too
large for operator-level analysis.  We select 20 counters according to three criteria:
(1)~coverage of every bottleneck axis in the roofline taxonomy (\secref{sec:roofline});
(2)~availability of a valid counter name on Ampere, Hopper, \emph{and} Blackwell (some counters
were renamed in the Blackwell GB202 silicon); (3)~non-redundancy (no two selected counters
correlate perfectly).

\subsection{The 20-Counter Set}

\sys collects 20 hardware counters covering all major bottleneck axes: memory bandwidth,
compute throughput, Tensor Core utilization, occupancy, latency, and register pressure.
The full counter list with bottleneck axis, aggregation method, and architecture-specific
name fallbacks is in \tabref{tab:counters} (\secref{sec:appendix:counters}).  Seven counters
require name fallbacks between Ampere, Hopper, and Blackwell; \sys queries all variants and
uses the first that returns a value.

\subsection{Duration-Weighted Aggregation}
\label{sec:aggregation}

An operator may dispatch multiple kernels with heterogeneous durations.  For rate and utilization
metrics (rows marked DW-mean in \tabref{tab:counters}), a naive arithmetic mean would give a
short-lived auxiliary kernel the same weight as the long-running compute kernel.  \sys uses
duration-weighted aggregation instead:
\begin{equation}
  \overline{m}_\text{op} = \frac{\displaystyle\sum_{i} m_i \cdot d_i}{\displaystyle\sum_{i} d_i}
  \label{eq:dwmean}
\end{equation}
where $m_i$ and $d_i$ are the metric value and duration of the $i$-th kernel attributed to the
operator.  This ensures that the dominant kernel---the primary optimization target---contributes
proportionally to the operator-level summary.  For additive quantities (bytes, instruction counts,
spills), \sys uses plain summation.  For per-kernel constants (register count, shared memory),
\sys takes the maximum across all kernels, since the highest-pressure kernel limits occupancy.

% -------------------------------------------------------
\section{Applying the Profile: Optimization Case Study}
\label{sec:pgo}

To demonstrate that the attributed hardware profiles produced in \secref{sec:attribution}
are actionable, we implement a profile-guided optimization workflow in which per-operator
counter values from \texttt{profile.json} drive the selection of FX graph transformations.
This section describes the optimization infrastructure and patterns; their empirical results
are in \secref{sec:experiments}.  The optimization layer is one possible downstream
application of the attribution pipeline---others include hardware-aware neural architecture
search, performance regression detection in CI, and kernel-level profiling of serving stacks
such as vLLM or SGLang.

\subsection{FX Pass Infrastructure}

\sys provides two composable components for graph-level optimization.
\textbf{UniqueSubgraphRegistry} partitions the FX graph by detected layer
structure, identifies structurally identical subgraphs via hash comparison,
and compiles a single representative per equivalence class.  Duplicate
partitions share the compiled callable, reducing compilation and ncu replay
time by $O(\text{num\_duplicate\_layers})$.

\textbf{FxPassRunner} applies user-defined pattern rewrites to unique
representatives only, then propagates the transformed compiled callables to
all structural duplicates:

\begin{lstlisting}[style=pystyle]
runner = FxPassRunner(registry)
for pattern_fn, replacement_fn in get_passes():
    runner.apply_pass(pattern_fn, replacement_fn)
\end{lstlisting}

Custom backends are registered via Torch Dynamo's backend hook and receive
the Aten IR \texttt{GraphModule} before Inductor lowering.

\subsection{Three-Level FX Pass Architecture}
\label{sec:threelevel}

Profile-guided optimizations are routed across three IR levels depending on the
graph representation at which each transformation is expressible.

\textbf{Level 1 --- Functional (pre-AOTAutograd).}
Fusion and operator substitution passes run here.  The key invariant is that the
Dynamo graph preserves the single shared-activation node for operations like a QKV
projection: at this level, three linear layers reading from one \texttt{LayerNorm}
output form a three-GEMM pattern that can be fused into one wide GEMM.
AOTAutograd decomposition shatters that shared node into per-consumer
\texttt{aten.view} copies, making the pattern undetectable at the Aten level.
SDPA canonicalization similarly requires the high-level \texttt{F.scaled\_dot\_product\_attention}
call node, which is replaced by Inductor-specific primitives after decomposition.

\textbf{Level 2 --- Aten IR (\texttt{\_aten\_inner\_compile}).}
Dtype casts, BatchNorm folding, and channels-last annotation run against the fully
decomposed \texttt{torch.ops.aten.*} node set.  An important subtlety: the
\texttt{example\_inputs} argument of \texttt{inner\_compile} may be \texttt{FakeTensor}
objects (symbolic values without storage), making it impossible to read weight values
directly.  Weight-reading passes (BN fold, pre-transposition) receive \texttt{real\_inputs}
threaded through \texttt{functools.partial} at the outer backend call site, before any
compilation begins.

\textbf{Level 3 --- Inductor config (\texttt{config\_patches}).}
Non-graph directives---\texttt{freezing}, \texttt{max\_autotune}, \texttt{max\_autotune\_gemm}---
are passed as a scoped configuration dictionary to \texttt{compile\_fx} rather than
mutating the global \texttt{torch.\_inductor.config}.  This ensures the policy applies
only to the specific compilation and cannot leak across compilation units.

The three levels compose as a fixed funnel:
\begin{lstlisting}[style=pystyle]
@register_backend
def my_backend(gm, example_inputs):
    gm = _run_functional_passes(gm)               # Level 1
    inner = functools.partial(_aten_inner_compile,
                              real_inputs=list(example_inputs))
    return compile_fx(gm, example_inputs,          # Level 2 + 3
                      inner_compile=inner,
                      config_patches={"freezing": True})
\end{lstlisting}

Each pass degrades gracefully: a \texttt{try}/\texttt{except} wrapper logs
\texttt{WARNING} and returns the graph unchanged if the target pattern is absent,
so partial optimization is always preferred over a compilation failure.
\texttt{gm.graph.lint()} is called after each mutation to validate graph invariants.
When \texttt{UniqueSubgraphRegistry} detects repeated layer structure, the funnel
runs once per unique representative and shares the compiled callable across all
structural duplicates---each FX pass is therefore authored once for the full model.

\subsection{Optimization Patterns}
\label{sec:passes}

Hardware-evidence-motivated optimization patterns are implemented per-workload
in each example's \texttt{*\_optimized.py}.  Patterns fall into two categories:
\emph{pre-compile model modifications} applied before \texttt{torch.compile} is
called, and \emph{FX graph pattern rewrites} applied inside a custom backend via
\texttt{torch.fx.subgraph\_rewriter.replace\_pattern}.  Each pattern fires only
when the profiler has diagnosed the specific counter symptom that motivates it.
\tabref{tab:passes} summarizes the patterns demonstrated across the evaluated
workloads.

\begin{table}[t]
\centering
\caption{Hardware-evidence-motivated optimization patterns demonstrated across
  the evaluated workloads.  ``Pre-compile'' patterns modify the \texttt{nn.Module}
  or tensor before \texttt{torch.compile}; ``FX rewrite'' patterns modify the
  Aten IR graph inside a custom backend.}
\label{tab:passes}
\small
\setlength{\tabcolsep}{4pt}
\begin{tabular}{p{2.0cm}p{1.5cm}p{2.8cm}p{3.3cm}}
\toprule
\textbf{Pattern} & \textbf{Type} & \textbf{Counter Evidence} & \textbf{Transformation} \\
\midrule
BF16 promotion
  & Pre-compile
  & \texttt{tensor\_core\_active\_pct < 5\%}
  & \texttt{model.to(bfloat16)} before compile; routes GEMM to HMMA Tensor Core path \\[4pt]
\texttt{channels\_last} layout
  & Pre-compile
  & \texttt{convertTensor\_kernel} in \nsys trace
  & \texttt{model.to(channels\_last)}; eliminates cuDNN NCHW$\to$NHWC coercion \\[4pt]
BatchNorm folding
  & Pre-compile
  & \texttt{\_native\_batch\_norm\_legit\_no\_training} in profile
  & Fold BN running statistics into preceding Conv2d weights; replace module with \texttt{Identity} \\[4pt]
Weight pretranspose
  & FX rewrite
  & Low \texttt{sm\_throughput\_pct} on TN GEMM
  & Replace \texttt{aten.t(W)} + \texttt{aten.mm} with pre-computed \texttt{W$^\top$.contiguous()} buffer \\[4pt]
max-autotune
  & Compiler setting
  & Sub-optimal tile selection (SM throughput below peak)
  & \texttt{config\_patches=\{`max\_autotune': True\}} in \texttt{compile\_fx} \\
\bottomrule
\end{tabular}
\end{table}


\subsection{Closed-Loop Validation Protocol}
\label{sec:validation}

The full optimization cycle consists of four steps:
\begin{enumerate}
  \item \textbf{Profile}: capture baseline \texttt{profile.json} via \nsys + \ncu pipeline. The \nsys capture runs with GPU clocks locked to the probe-and-cache frequency (\secref{sec:clocklocking}); the same cached clock is reused for the optimized re-capture in step~4, ensuring duration ratios reflect architectural changes rather than thermal variation.
  \item \textbf{Diagnose}: \texttt{profile.json} is examined per-operator; per-operator hardware
    counter values are used to produce ranked optimization proposals (\texttt{optimizations.json})
    with evidence citations and dependency ordering.
  \item \textbf{Transform}: implement a workload-specific optimized backend and
    run \texttt{torch.compile(model, backend="<workload>\_opt")} to apply the patterns and compile.
  \item \textbf{Validate}: capture optimized \texttt{profile\_optimized.json} and generate a delta
    report comparing per-operator counter values, flagging any regression.
\end{enumerate}

This protocol makes every optimization auditable: the delta report shows not just speedup, but
\emph{which counter improved}, providing evidence that the diagnosed bottleneck was actually
eliminated rather than merely shifted to a different operator.

% -------------------------------------------------------
\section{Design Challenges}
\label{sec:challenges}

Implementing \sys required resolving several non-obvious engineering obstacles.
We document them here as they motivate the design decisions in \secref{sec:attribution}
and \secref{sec:pgo}, and because the failure modes are not documented elsewhere.

\subsection*{CUPTI Exclusivity: Two-Phase Capture}

\nsys and \texttt{torch.profiler} both consume CUPTI for kernel timing and cannot run simultaneously within the same process.  An early design ran the torch.profiler correlation pass \emph{inside} an \texttt{nsys profile} invocation, expecting both to collect independently.  In practice, CUPTI preempts the \texttt{torch.profiler} callback with no error; the correlation pass produces a \texttt{.corr.json} with zero entries, and the silence makes the failure invisible.  The pipeline was redesigned as two sequential invocations: Phase~1 is a plain Python run (no nsys) that captures torch.profiler output and writes \texttt{.corr.json}; Phase~2 runs under \texttt{nsys profile} and reuses the Inductor cache from Phase~1.

\subsection*{NVTX v3 Static Linking: ncu Range Mode Abandoned}

An appealing approach was to use \ncu's \texttt{--filter-by-nvtx-range} to replay only the kernels within each \texttt{aten::} NVTX range, collecting per-operator hardware counters without a full-workload replay.  PyTorch~2.x statically links NVTX~v3 inside \texttt{libtorch\_cuda.so}, whereas \ncu's injection hooks target the dynamically-loaded NVTX library.  As a result, \ncu's injection mechanism never observes any NVTX events; range filters silently capture nothing regardless of \texttt{emit\_nvtx} activity.  Application-mode replay (\texttt{--replay-mode application}) is the only path that reliably collects counters for Inductor-compiled workloads.

\subsection*{Functional vs.\ Aten IR Level Mismatch}

The initial backend implementation applied all FX passes at a single IR level.  Profiling the GPT-2 workload revealed that QKV fusion patterns were not found at runtime: after AOTAutograd decomposition, the shared \texttt{q/k/v} activation node had been replicated into per-consumer \texttt{aten.view} copies, destroying the three-GEMM pattern the fusion matcher expected.  This motivated the three-level architecture (\secref{sec:threelevel}).  The fix was non-trivial because the Dynamo functional level and the post-AOTAutograd Aten level look superficially similar; the difference only surfaces at profiling time when a supposedly-applied pass produces no measurable counter change.

\subsection*{FakeTensor Opacity in \texttt{\_aten\_inner\_compile}}

The \texttt{example\_inputs} argument of Inductor's \texttt{inner\_compile} hook may be \texttt{FakeTensor} objects---symbolic values that carry shape and dtype metadata but have no allocated storage.  An initial implementation of BatchNorm folding attempted to read \texttt{.data} from these tensors and crashed with storage-access errors.  The fix threads \texttt{real\_inputs = list(example\_inputs)} from the outer backend entry point before any \texttt{compile\_fx} call, so weight-reading passes operate on concrete tensors while Inductor's trace receives the \texttt{FakeTensor} inputs it expects.

\subsection*{GPU Clock Drift}

Initial example profiles were captured without clock management.  Baseline and optimized captures made at different times---with the optimized workload running cooler due to fewer scalar SIMT FLOPs---showed spurious duration differences on operators that were architecturally unchanged.  In one case the apparent speedup on an unmodified kernel exceeded 1.2$\times$, solely from clock drift.  The probe-and-lock mechanism (\secref{sec:clocklocking}) was added after this was identified as a systematic bias in the before/after comparison.

\subsection*{NVTX Warmup Attribution Contamination}

The layer deduplication NVTX wrapping initially emitted \texttt{layer::unique::} and \texttt{layer::duplicate::} ranges for every forward pass, including warmup iterations.  \texttt{ManifestBuilder}'s \texttt{\_tag\_layer\_partitions} attributed kernels launched during warmup to their enclosing layer-partition ranges, introducing warmup-only compilation artifacts into operator profiles and inflating attribution coverage.  The fix gates range emission behind an \texttt{\_NVTX\_STATE[`active']} flag set only after warmup completes.  Both bugs (clock drift and warmup contamination) were discovered concurrently and corrected in the same pipeline revision.

\subsection*{Kernel Invocation Count Invariant}

Invocation-order matching (\algref{alg:iom}) is correct only when the \texttt{--warmup-iters} and \texttt{--measure-iters} counts are identical between the \nsys capture and the \ncu replay.  A mismatch shifts kernel invocation indices, causing hardware counters to be assigned to the wrong operator call without any runtime error.  This invariant was violated during early development when a second \nsys capture used different iteration counts without regenerating the manifest; the resulting \texttt{profile.json} contained plausible-looking but incorrect counter values.  The pipeline now validates iteration counts before initiating replay.

\subsection*{sudo Environment Stripping Under \texttt{env\_reset}}

\ncu requires elevated privileges to access GPU performance counters.  Running \texttt{sudo ncu} on a system with the default \texttt{sudoers env\_reset} rule strips \texttt{PYTHONPATH} from the subprocess environment; the workload script imports from \texttt{nvidia.operator\_profiler}, so the subprocess exits with \texttt{ModuleNotFoundError} with no counter output and no obvious connection to the import path.  The fix uses \texttt{sudo env KEY=VAL \ldots\ ncu} to forward environment variables explicitly, bypassing the \texttt{env\_keep} list entirely.

\section{Evaluation: Attribution Quality and Optimization Case Studies}
\label{sec:experiments}

We evaluate \sys on two axes: (1) \emph{coverage}---what fraction of kernel runtime is
correctly attributed to \texttt{aten::} operators at each confidence tier; and
(2) \emph{utility}---whether attributed profiles provide sufficient signal to ground and
verify optimization decisions. 

\subsection{Experimental Setup}

\textbf{Hardware.}  All experiments run on a single NVIDIA RTX PRO 6000 Blackwell GPU; full hardware specifications are in \secref{sec:appendix:hwspec}.

\textbf{Workloads.}  \tabref{tab:workloads} describes the three workloads evaluated.  Each
isolates a distinct operator category and hardware challenge.

\begin{table}[h]
\centering
\caption{Case study workloads.}
\label{tab:workloads}
\small
\begin{tabular}{@{}lcp{3.0cm}p{5.2cm}@{}}
\toprule
\textbf{Workload} & \textbf{Batch} & \textbf{Dimensions} & \textbf{Profile Pattern} \\
\midrule
GPT-2 (12-layer)   & 4  & hidden=768, FFN=3072, 12 heads, seq=128 & TC idle (\texttt{tensor\_core\_active\_pct}$\approx$0\%) on all GEMM families \\
SDPA Attention     & 8  & seq=512, dim=512, 8 heads & Low occupancy (16.6\%) + TC idle GEMMs + FP32 attention kernel \\
LSTM Seq.\ Encoder & 32 & seq=128, hidden=512, 2 layers & Structural: Inductor LSTM decomposition into 1,280 scalar SIMT GEMMs \\
\bottomrule
\end{tabular}
\end{table}

\textbf{Measurement protocol.}  Per-operator durations are \nsys CUPTI GPU-domain kernel times.
Each workload runs 2 warm-up iterations followed by 2 measured iterations; reported times are
the sum over measured iterations, and speedup ratios are point estimates without variance
reporting.  At locked clocks, kernel scheduling jitter for these workloads is ${<}$2\% across
repeated runs; point estimates suffice to distinguish the 1.76$\times$--3.63$\times$ speedups
reported here from the 1.0$\times$ null case.  GPU clocks are
locked to the probe-and-cache frequency (\secref{sec:clocklocking}) for the duration of both
the baseline and optimized \nsys captures; both captures lock to the same frequency, so
duration ratios are clock-immune.  \ncu replay is performed in application mode
(\texttt{--replay-mode application}), collecting all 20 metrics in 4--8 passes; \ncu
replay durations are 2--5$\times$ elevated relative to real wall-clock execution and are not
used for speedup reporting.  Layer deduplication is active for models with repeated structure
(\eg GPT-2), profiling unique representatives and propagating metrics to duplicates.

\textbf{Speedup scope.}  All speedup figures in \tabref{tab:speedup} are computed over
\emph{attributed kernel wall time}: the sum of CUPTI GPU-domain durations for kernels that
received an attribution label.  Unattributed kernels (e.g., opaque cuDNN RNN internals in the
LSTM baseline), CPU dispatch overhead, and inter-operator synchronization are excluded.  For
workloads with high unattributed fractions (LSTM: $>$85\% baseline), the reported speedup
reflects the attributed portion only and may differ from end-to-end wall-clock improvement.

\textbf{Baselines.}  Each workload is compared as:
\begin{enumerate*}[label=(\arabic*)]
  \item \emph{Inductor FP32 baseline}: \texttt{torch.compile(backend="inductor")} with
    default settings, FP32 precision;
  \item \emph{Optimized}: workload-specific backend applying BF16 promotion and other
    hardware-evidence-motivated patterns.
\end{enumerate*}
Speedups are reported as \emph{Inductor FP32 baseline} $\div$ \emph{Optimized} in attributed
kernel wall time at locked clocks.

\subsection{Attribution Coverage}
\label{sec:coverage}

Attribution coverage for the three case study workloads is shown in \figref{fig:attribution}.
GPT-2 and SDPA Attention achieve near-complete attribution with the correlation pass active.
LSTM's $>$85\% unattributed fraction is the expected result for that dispatch path; the cause
and corrective strategy are examined in \secref{sec:cs:lstm}.

\begin{figure}[t]
\centering
\begin{tikzpicture}
\begin{axis}[
  xbar stacked,
  width=0.88\columnwidth,
  height=3.6cm,
  bar width=13pt,
  xmin=0, xmax=100,
  xlabel={Kernel runtime attributed (\%)},
  xlabel style={font=\small},
  ytick=data,
  yticklabels={
    LSTM Seq.\ Enc.,
    SDPA Attention,
    GPT-2
  },
  yticklabel style={font=\footnotesize},
  tick label style={font=\footnotesize},
  axis x line=bottom,
  axis y line=left,
  xmajorgrids=true,
  grid style={gray!20},
  legend style={
    at={(0.5,-0.28)}, anchor=north,
    font=\scriptsize, draw=gray!50,
    fill=white, fill opacity=0.9,
    legend columns=3,
  },
]

\addplot[fill=highcol, draw=highcol!80!black] coordinates {
  (4,   0)   % LSTM (opaque cuDNN RNN)
  (35,  1)   % SDPA
  (40,  2)   % GPT-2
};

\addplot[fill=medcol!70, draw=medcol!80!black] coordinates {
  (8,   0)   % LSTM (compiled head fraction)
  (62,  1)   % SDPA
  (57,  2)   % GPT-2
};

\addplot[fill=ucol!80, draw=gray!60] coordinates {
  (88,  0)   % LSTM --- >85% cuDNN RNN opacity
  (3,   1)   % SDPA
  (3,   2)   % GPT-2
};

\legend{HIGH (torch.profiler IOM), MEDIUM (NVTX / Inductor), UNATTRIBUTED}

\end{axis}
\end{tikzpicture}
\caption{\textbf{Attribution coverage for the three case study workloads.}  GPT-2 and SDPA
  Attention achieve $>$95\% attributed runtime with the correlation pass active.  LSTM's
  $>$85\% UNATTRIBUTED fraction correctly reflects opaque cuDNN RNN dispatch
  (see \secref{sec:cs:lstm}).  Inductor fusion enrichment (\secref{sec:heuristic})
  typically adds 5--15 percentage points for compiled workloads.}
\label{fig:attribution}
\end{figure}

\subsection{Profiling Overhead}

\tabref{tab:overhead} characterizes the overhead of the two-stage profiling pipeline, measured
for GPT-2 on RTX PRO 6000 Blackwell.  Layer deduplication (12 transformer blocks $\to$ 1 unique representative)
provides a 12$\times$ reduction in \ncu replay time for this workload.

\begin{table}[h]
\centering
\caption{Two-stage profiling overhead for GPT-2 (RTX PRO 6000 Blackwell, B=4, seq=128, 2 measure
  iterations, locked clocks) with layer deduplication active.  Profiling is an offline, one-time cost.}
\label{tab:overhead}
\small
\begin{tabular}{llll}
\toprule
\textbf{Phase} & \textbf{Wall Time} & \textbf{Overhead} & \textbf{Note} \\
\midrule
Clock probe (first run only)     & ${\sim}$5\,s       & one-time           & Probe-and-cache; reused across all captures \\
\nsys capture (2 iters)          & ${\sim}$10\,s      & $3\times$ baseline & NVTX tracing adds ${\sim}$2$\times$ overhead \\
\ncu replay (unique layers only) & ${\sim}$3--5\,min  & serializes kernels & 12$\times$ speedup vs.\ no dedup (\S\ref{sec:dedup}) \\
Attribution \& aggregation (CPU) & ${<}$5\,s          & negligible         & SQL query + interval tree build \\
\textbf{Total time to profile}   & ${\sim}$4--6\,min  & one-time cost      & Amortized over iterative optimization \\
\bottomrule
\end{tabular}
\end{table}

The profiling cost is commensurate with an offline, iterative optimization workflow: a
single profiling run produces the attributed profile; the optimized backend is implemented
and validated against a second run at matched clock frequencies.  Re-profiling is required
only when the operator set changes between model revisions.

\subsection{Optimization Case Studies}
\label{sec:validation}

The following case studies demonstrate how \texttt{profile.json} grounds optimization
decisions.  Each study identifies the counter symptom read from the attributed profile,
the transformation selected in response, and the before/after counter delta that confirms
bottleneck elimination.  This renders every optimization auditable: the delta report records
\emph{which counter improved}, confirming the diagnosed bottleneck was eliminated rather than
shifted to a different operator.  \tabref{tab:speedup} summarizes attributed kernel wall time
speedups at locked GPU clocks.

\begin{table}[h]
\centering
\caption{Forward-pass attributed kernel wall time and speedup on NVIDIA RTX PRO 6000 Blackwell
 .  All times are from \nsys CUPTI records at locked GPU clocks (not \ncu replay).
  Speedup = Inductor FP32 $\div$ Optimized (attributed kernel time only; see \S\ref{sec:experiments}
  for scope).  Each result is a point estimate from 2 measured iterations.}
\label{tab:speedup}
\small
\setlength{\tabcolsep}{5pt}
\begin{tabular}{lrrr}
\toprule
\textbf{Workload} & \textbf{Inductor FP32 Baseline} & \textbf{Optimized}
                  & \textbf{Speedup} \\
\midrule
GPT-2 (12-layer)        & 7,313\,$\mu$s  & 4,161\,$\mu$s  & \textbf{1.76$\times$} \\
SDPA Attention          & 906\,$\mu$s    & 405\,$\mu$s    & \textbf{2.24$\times$} \\
LSTM Seq.\ Encoder      & 14,531\,$\mu$s & 3,998\,$\mu$s  & \textbf{3.63$\times$} \\
\bottomrule
\end{tabular}
\end{table}

\tabref{tab:hwcounters} reports the key hardware counter deltas that motivated and confirmed
each optimization; each case study below identifies which counter symptom drove the decision.

\begin{table}[h]
\centering
\caption{Key hardware counter deltas for all three case study workloads
  (RTX PRO 6000 Blackwell, \ncu application-mode replay).
  GEMM counters are duration-weighted aggregates (\equref{eq:dwmean}).}
\label{tab:hwcounters}
\small
\setlength{\tabcolsep}{4pt}
\begin{tabular}{llrrl}
\toprule
\textbf{Workload} & \textbf{Counter} & \textbf{Baseline} & \textbf{Optimized}
  & \textbf{Bottleneck Resolved} \\
\midrule
\multicolumn{5}{l}{\textit{GPT-2 (B=4, seq=128)}} \\
& TC active, GEMMs (\%)          & 0.0\%        & 40.5\%       & FP32 SIMT $\to$ BF16 Tensor Core \\
& Registers/thread (GEMMs)       & 210          & 96           & BF16 tiles need fewer registers \\
& Achieved occupancy, GEMMs (\%) & 16.5\%       & ${\sim}$30\% & Register pressure relief \\
& SM throughput, GEMMs (\%)      & ${\sim}$40\% & ${\sim}$65\% & Correct compute pipeline \\
& Attention kernel                & \texttt{fmha\_f32} & \texttt{fmha\_f32} & Stays FP32; SDPA pass removes mask prep \\
\midrule
\multicolumn{5}{l}{\textit{SDPA Attention (B=8, T=512, D=512)}} \\
& TC active, QKV GEMMs (\%)      & 0.0\%  & 19--21\% & FP32 SIMT $\to$ BF16 Tensor Core \\
& Occupancy, QKV GEMMs (\%)      & 16.6\% & 73.9\%   & 3 serial $\to$ 1 wide launch (QKV fusion) \\
& Attention kernel                & \texttt{fmha\_f32} & \texttt{fmha\_bf16} & BF16 promotion activates bf16 fmha \\
\midrule
\multicolumn{5}{l}{\textit{LSTM Seq.\ Encoder (B=32, seq=128, hidden=512)}} \\
& TC active, recurrent GEMMs (\%)  & 0.0\%   & 23.8\% & Scalar SIMT $\to$ cuDNN fused RNN \\
& Occupancy, recurrent GEMMs (\%)  & 9.4\%   & 39.4\% & M=32 $\to$ M=4096 (batch timesteps) \\
& \texttt{splitKreduce} epilogues  & 1{,}280 & 0      & Eliminated by cuDNN routing \\
\bottomrule
\end{tabular}
\end{table}

\subsection{Case Study 1 --- GPT-2: Tensor Core Idle $\to$ BF16 Promotion}
\label{sec:cs:gpt2}

\tabref{tab:transformer} reports per-operator attributed time for GPT-2 small (12 transformer
blocks, hidden=768, 12 heads, FFN=3072, B=4, seq=128) on the RTX PRO 6000 Blackwell.
All times are \nsys-derived GPU kernel wall times at locked clocks (1845\,MHz graphics).
The baseline is dominated by 96 \texttt{cutlass\_80\_simt\_sgemm} kernels; the optimized
backend replaces them with \texttt{cutlass\_80\_tensorop\_bf16} and
\texttt{cutlass\_80\_wmma\_tensorop\_bf16} kernels (\tabref{tab:hwcounters}).

\begin{table}[h]
\centering
\caption{Per-operator-class attributed time for GPT-2 (RTX PRO 6000 Blackwell, B=4, seq=128,
  locked clocks).  Baseline: Inductor FP32 (built-in dedup backend).  Optimized: BF16 dtype
  promotion + SDPA canonicalization + Inductor freezing (\texttt{gpt2\_opt} backend).
  ``Triton casts'' are fp32$\leftrightarrow$bf16 cast kernels introduced by the bf16 promotion
  that the cast-cancellation pass (OPT-4) was unable to eliminate.}
\label{tab:transformer}
\small
\begin{tabular*}{\textwidth}{@{\extracolsep{\fill}}lrrr}
\toprule
\textbf{Operator class} & \textbf{Baseline} & \textbf{Optimized} & \textbf{Speedup} \\
\midrule
GEMM family (\texttt{cutlass\_80\_simt\_sgemm}, 96 kernels) & 6,401\,$\mu$s & 2,695\,$\mu$s & \textbf{2.37$\times$} \\
Attention (\texttt{fmha\_cutlassF\_f32})                    &   485\,$\mu$s &   408\,$\mu$s & 1.19$\times$ \\
LayerNorm                                                   &   149\,$\mu$s &   120\,$\mu$s & 1.24$\times$ \\
Triton casts + elementwise + epilogues                      &   279\,$\mu$s &   938\,$\mu$s & new overhead \\
\midrule
\textbf{Total}                                              & \textbf{7,313\,$\mu$s} & \textbf{4,161\,$\mu$s} & \textbf{1.76$\times$} \\
\bottomrule
\end{tabular*}
\end{table}

The attributed breakdown makes the optimization target unambiguous: 87.5\% of total attributed
runtime (6{,}401 of 7{,}313\,$\mu$s) resides in the GEMM family.  An aggregate
workload-level \texttt{tensor\_core\_active\_pct} reading would confirm that Tensor Cores are
idle, but would not identify whether that idle time is concentrated in GEMMs or distributed
across attention, LayerNorm, and epilogue kernels.  The per-operator breakdown enables the
optimizer to predict, before writing any backend code, that BF16 promotion targeting the GEMM
family will recover a near-proportional fraction of total runtime.

The primary overhead introduced by the optimization is the fp32$\leftrightarrow$bf16 cast
traffic shown in the ``Triton casts'' row: the BF16 promotion pass (\texttt{OPT-1}) inserts
casts around each GEMM node, and the cast-cancellation pass (\texttt{OPT-4}) failed to
eliminate them because \texttt{OPT-1} restores the output dtype directly rather than emitting
per-node inverse cast pairs.  This expanded the Triton elementwise kernel family from 72 to
288 invocations, adding 0.66\,ms and reducing the net attributed-kernel gain from an estimated
${\sim}$2.1$\times$ to the observed 1.76$\times$.

\begin{figure}[t]
\centering
\begin{tikzpicture}
\begin{axis}[
  xbar,
  width=0.88\columnwidth,
  height=4.8cm,
  bar width=14pt,
  xmin=0, xmax=3.2,
  xlabel={Speedup ($\times$, nsys locked-clock wall time)},
  xlabel style={font=\small},
  ytick=data,
  yticklabels={
    {Remaining (LayerNorm + epilogue)},
    {Attention (fmha)},
    {GEMM family}
  },
  yticklabel style={font=\footnotesize},
  tick label style={font=\footnotesize},
  nodes near coords,
  nodes near coords style={font=\scriptsize, anchor=west},
  nodes near coords align={horizontal},
  every node near coord/.append style={xshift=2pt},
  axis x line=bottom,
  axis y line=left,
  xmajorgrids=true,
  grid style={gray!20},
  extra x ticks={1.76},
  extra x tick labels={Total 1.76$\times$},
  extra x tick style={
    grid=major,
    grid style={red!60!black, dashed, thick},
    tick label style={font=\scriptsize, text=red!70!black, rotate=90, anchor=west},
  },
  legend style={
    at={(0.97,0.03)}, anchor=south east,
    font=\scriptsize, draw=gray!50,
    fill=white, fill opacity=0.9,
  },
]

\addplot[fill=orange!70, draw=orange!80!black, bar shift=0pt] coordinates { (2.37, 2) };
\addlegendentry{BF16 GEMM (TC 0\%$\to$40.5\%)}

\addplot[fill=teal!50, draw=teal!70!black, bar shift=0pt] coordinates { (1.19, 1) };
\addlegendentry{BF16 attention}

\addplot[fill=blue!45, draw=blue!60!black, bar shift=0pt] coordinates { (1.22, 0) };
\addlegendentry{BF16 minor gain}

\end{axis}
\end{tikzpicture}
\caption{\textbf{Per-operator-class speedup for GPT-2} (RTX PRO 6000 Blackwell, B=4, seq=128,
  locked-clock nsys wall time).
  \textcolor{orange!80!black}{Orange}: GEMM family (2.37$\times$).
  \textcolor{teal!70!black}{Teal}: attention (1.19$\times$).
  \textcolor{blue!60!black}{Blue}: remaining operators (1.22$\times$).
  Dashed red line: overall \textbf{1.76$\times$}.
  The gap between the GEMM gain and the end-to-end total is explained by new
  fp32$\leftrightarrow$bf16 cast overhead (\tabref{tab:transformer}).}
\label{fig:speedup_bars}
\end{figure}

The attention kernel remains FP32 (\texttt{fmha\_cutlassF\_f32}); the attention speedup
derives from removing the causal-mask construction subgraph (SDPA canonicalization pass),
not a kernel dtype change.  Unlike on A100, the Blackwell attention kernel did not switch to
a lower-register BF16 path---the dominant lever on this workload is exclusively the GEMM
family (\tabref{tab:hwcounters}).

\subsubsection*{Diagnosis and Transformation}

The \texttt{gpt2\_opt} backend applies four optimization passes; the three that produced
measurable counter improvements are described below.  A pass is credited with an improvement
only when its \texttt{status} field is \texttt{APPLIED}, the target counter changes in the
expected direction, and the corresponding operator exhibits reduced runtime.

\textbf{BF16 dtype promotion (OPT-1, +2.37$\times$ on the GEMM family):}
Casting \texttt{mm}/\texttt{addmm} operands to bf16 reroutes all four projection GEMM families
off the FP32 SIMT path and onto the Blackwell BF16 Tensor Core path; hardware counter
confirmation is in \tabref{tab:hwcounters}.

\textbf{SDPA canonicalization (OPT-2):}
Switching \texttt{F.scaled\_dot\_product\_attention} to \texttt{is\_causal=True} removes the
48-launch memory-bound mask-construction subgraph.  This pass runs at the functional level
(\secref{sec:threelevel}) because the high-level SDPA call node is only present before
AOTAutograd decomposition.

\textbf{Inductor freezing (OPT-3):}
\texttt{freezing=True} constant-folds LayerNorm affine weights and projection biases, letting
CUTLASS select fused-epilogue GEMM templates.  The standalone fused-bias Triton family (46
launches, 0.17\,ms) is folded into the GEMM epilogue path in the optimized profile.

\textbf{Cast cancellation (OPT-4, NOT\_APPLIED):}
The cast-cancellation pass detected no per-node inverse cast pairs to cancel because OPT-1
restores the output dtype directly.  The residual fp32$\leftrightarrow$bf16 cast traffic ($+$0.66\,ms,
288 new Triton kernels) is the dominant residual opportunity---revising OPT-1 to maintain
activations in bf16 across consecutive GEMMs would recover most of this overhead, increasing
the net attributed-kernel speedup from the observed 1.76$\times$ toward an estimated
${\sim}$2.1$\times$.

\subsection{Case Study 2 --- SDPA Attention: Low Occupancy and QKV Fusion}
\label{sec:cs:sdpa}

The SDPA workload (multi-head attention, B=8, T=512, D=512, 8 heads) shows a 2.24$\times$
attributed-kernel speedup (\tabref{tab:speedup}).  The baseline concentrates 96.7\% of
attributed time in two operator families: eight \texttt{aten::mm} GEMMs (66.2\%, all TC-idle)
and two SDPA calls (30.5\%, FP32 \texttt{fmha\_cutlassF\_f32\_aligned\_64x64\_sm80}).

A workload-level occupancy reading would reveal low average occupancy but cannot identify
which operator is responsible or why.  The attributed profile isolates occupancy\,=\,16.6\%
to the Q/K/V projection GEMMs specifically, and the FX graph reveals that all three
projections read from a single shared \texttt{LayerNorm} activation---the structural
precondition for QKV fusion.  This causal chain (per-operator counter $\to$ graph-structural
diagnosis $\to$ specific FX pass) is not available from aggregate profiling alone.

Three optimization passes were successfully applied.  \textbf{BF16 dtype promotion} (OPT-1)
rerouted the GEMMs off the FP32 SIMT path and switched the SDPA kernel from
\texttt{fmha\_cutlassF\_f32} to \texttt{fmha\_cutlassF\_bf16}
(${\sim}$\textbf{3.0$\times$} on the attention kernel; counter deltas in \tabref{tab:hwcounters}).
\textbf{QKV fusion} (OPT-2) merged the three Q/K/V projections into a single
\texttt{[512, 1536]} GEMM (three serial 128-block launches $\to$ one 384-block launch),
raising achieved occupancy to 73.9\% (\tabref{tab:hwcounters}).
\textbf{Inductor freezing} (OPT-3) arranged the projection weights in aligned memory
order, enabling the bf16 GEMM autotuner to select aligned templates.

The total attributed-kernel speedup (2.24$\times$) is bounded by newly-introduced auxiliary
kernels---bf16 cast boundaries, the QKV projection split-and-concatenation, and a
pre-transposed frozen-weight conversion---which collectively account for 9.1\% of the
optimized profile runtime.

Unlike the A100 result for FMHA (where register spills were the dominant SDPA bottleneck on
A100), on Blackwell the primary mechanism is Tensor Core activation on the GEMMs and the
bf16 attention kernel, not register pressure reduction.

\subsection{Case Study 3 --- LSTM Sequence Encoder: Structural Miscompilation}
\label{sec:cs:lstm}

The LSTM workload (2-layer LSTM, B=32, seq=128, hidden=512) produces the largest speedup in
the suite (\tabref{tab:speedup}: 3.63$\times$), and uniquely illustrates the framework's
ability to diagnose a structural inefficiency that is undetectable by timing-only profilers.

Under \texttt{torch.compile(backend="inductor")}, \texttt{nn.LSTM} triggers a hard Dynamo
graph break (\texttt{graph\_count = 0}): the operator enters via \texttt{torch.\_VF.lstm}, a
C extension entry point that Dynamo cannot graph-capture.  Inductor therefore decomposes the
LSTM into a per-timestep unrolled loop---1,280 small-matrix FP32 GEMMs (M=32) on the scalar
SIMT path (\texttt{tensor\_core\_active\_pct = 0.0}, achieved occupancy 9.4\%, 1,280 matching
\texttt{cublasLt::splitKreduce\_kernel} epilogues).  Attribution is $>$85\% UNATTRIBUTED because
cuDNN's internal RNN kernels carry no NVTX annotations.

The corrective strategy is structural: given that \texttt{nn.LSTM} unconditionally forces a
Dynamo graph break, the custom backend permits the recurrent region to execute eagerly,
routing it through cuDNN's fused Tensor-Core RNN rather than Inductor's scalar decomposition
path.  cuDNN batches the input-to-hidden projection across all 128 timesteps in one
launch (M=4096 instead of M=32), achieving \textbf{3.73$\times$} on the recurrent GEMMs and
eliminating all 1,280 \texttt{splitKreduce} epilogues, which are absent from the optimized
profile and reduce attributed execution time by 2.26\,ms; hardware counter deltas are in
\tabref{tab:hwcounters}.

This workload illustrates a scenario in which all FX graph passes report
\texttt{NOT\_APPLIED}---because the backend callback is never invoked for the
eagerly-dispatched LSTM region---yet the resulting speedup is the largest observed across
the evaluation suite.  The improvement is attributed to the structural routing decision
(eager dispatch to cuDNN), not to an FX graph transformation; the profile-guided workflow
correctly identified the bottleneck and the optimization follows directly from that
diagnosis.

This distinguishes the attributed pipeline from timing-only profiling.  A kernel-level timer
would expose 1{,}280 slow small-matrix GEMMs and motivate direct kernel optimization---tiling,
BF16 cast, epilogue fusion.  The attribution-tier signal ($>$85\% UNATTRIBUTED from a
cuDNN-opaque path) instead identifies the problem as structural: the LSTM body is not
reachable by any FX-level transformation, and the only valid fix is a dispatch routing
decision.  The attribution tier is the observable that distinguishes these two diagnoses.


\section{Conclusion}
\label{sec:conclusion}

We presented \sys, a hardware attribution pipeline that bridges the gap between PyTorch's
operator-level abstractions and NVIDIA's kernel-level hardware performance counters.  The
core contribution is a mapping from CUDA kernels to \texttt{aten::} operators via NVTX
temporal enclosure with per-stream interval trees, invocation-order matching across
incompatible clock domains, and Inductor debug-artifact enrichment for compiled models.
The result is \texttt{profile.json}: a hardware-attributed operator profile mapping each
\texttt{aten::} operator to its constituent kernels and their full hardware counter state,
produced automatically without manual kernel-to-operator correlation.  A curated 20-counter
metric set with duration-weighted aggregation and architecture-specific fallbacks covers all
hardware bottleneck axes; GPU clock locking ensures reproducible per-operator duration
measurements; and layer deduplication reduces \ncu replay time by
$O(\text{num\_duplicate\_layers})$ for models with repeated structure.

To demonstrate that attributed profiles are actionable, we implemented a profile-guided
FX pass framework (\secref{sec:pgo}) and applied it across three case studies on an NVIDIA
RTX PRO 6000 Blackwell Server Edition.  Each optimization decision was grounded in a
specific counter reading from the attributed profile, and each claimed improvement was
verified by re-profiling at identical locked clocks:
\textbf{1.76$\times$} on GPT-2 (BF16 GEMM promotion, \texttt{tensor\_core\_active\_pct}:
0\%$\to$40.5\%);
\textbf{2.24$\times$} on SDPA Attention (QKV fusion raises occupancy 16.6\%$\to$73.9\%;
SDPA kernel 3.0$\times$); and
\textbf{3.63$\times$} on LSTM Sequence Encoder (structural re-routing from Inductor's
decomposed SIMT path to cuDNN's fused Tensor-Core RNN, diagnosable only via the attribution
data).  The LSTM case is the clearest illustration of the pipeline's value: the bottleneck
is invisible to timing-only profilers but apparent from the high unattributed fraction and
scalar GEMM signatures in the attributed profile.

\paragraph{Limitations.}
Attribution coverage degrades for cuDNN-backed operators (\texttt{nn.LSTM},
\texttt{nn.MultiheadAttention} under cuDNN) where internal kernel names carry no NVTX
correspondence and no Inductor debug metadata is available; expected unattributed rates of
80--90\% make hardware-counter-driven optimization inapplicable, and bottleneck diagnosis must
fall back to kernel-name inference.  Invocation-order matching requires a deterministic
workload: models with conditional control flow or dynamic dispatch can violate the
\texttt{(kernel\_name, invocation\_index)} invariant without producing a shape error;
\sys validates invocation counts before issuing counter assignments and raises
\texttt{InvocationCountMismatchError} on violation.  \ncu serializes kernels for replay,
making concurrent multi-stream profiles require separate per-stream runs.  All speedup
figures are point estimates from 2 measured iterations at locked clocks; run-to-run variance
is not characterized.  Optimization patterns are currently user-defined per-workload;
automatic pattern selection from counter evidence is left as future work.

\paragraph{Future work.}
Online profiling via CUPTI Perfworks streaming would eliminate the separate \ncu pass, enabling
profiling during ordinary training without a dedicated replay run.  Multi-GPU attribution with
NCCL stream synchronization events would extend the pipeline to distributed workloads.
Counter-guided neural architecture search---using per-operator hardware utilization as a fitness
signal rather than validation accuracy alone---represents a longer-term direction for hardware-aware
model design.  Finally, extending the counter name mapping and pass library to AMD ROCm (\texttt{rocprof}
+ \texttt{rocm-smi}) would generalize \sys beyond NVIDIA hardware.

\newpage
\bibliographystyle{plainnat}
\bibliography{reference}
\newpage
\appendix
\onecolumn
\section{Reproduction and Additional Details}
\label{sec:appendix:exp}

\subsection{Hardware Specification}
\label{sec:appendix:hwspec}

All experiments run on a single GPU:
\begin{itemize}[leftmargin=1.5em]
  \item \textbf{NVIDIA RTX PRO 6000 Blackwell}: ${\sim}$188 SMs (GB202),
    96\,GB GDDR7 (${\sim}$1.8\,TB/s memory bandwidth),
    5th-generation Tensor Cores (HMMA BF16/FP8), Blackwell architecture (sm\_120).
    Note: the \texttt{warp\_cycles\_per\_instruction} counter is absent from the
    Blackwell counter set; \texttt{eligible\_\allowbreak cycles\_pct} serves as the
    latency-bound indicator instead.
\end{itemize}

Software requirements: NVIDIA GPU (Ampere minimum), \nsys $\ge$ 2024.6,
\ncu $\ge$ 2025.4.1, PyTorch 2.11+, and CUDA 12.8.
Reproduction scripts, \texttt{profile.json} pairs, and per-example counter comparisons are provided in the accompanying repository here: 
\url{https://github.com/yottalabsai/Profiler}

\subsection{Hardware Counter Reference}
\label{sec:appendix:counters}

\begin{table}[h]
\centering
\caption{The 20 hardware counters collected by \sys, their bottleneck axis, and
what they measure.  ``Agg.'' denotes the aggregation method
  applied when multiple kernels are attributed to one operator (see \secref{sec:aggregation}).
  $\dagger$~\texttt{warp\_cycles\_per\_instr} is absent from the Blackwell counter set (sm\_120);
  \texttt{eligible\_\allowbreak cycles\_pct} (counter~14) serves as the latency-bound indicator on the
  RTX PRO 6000 used for all experiments reported here.}
\label{tab:counters}
\scriptsize
\setlength{\tabcolsep}{4pt}
\begin{tabular}{clllc}
\toprule
\textbf{\#} & \textbf{Profile Field} & \textbf{Bottleneck Axis} & \textbf{Measures} & \textbf{Agg.}\\
\midrule
1  & \texttt{gpu\_time\_duration\_ns}      & Latency          & GPU wall time                & SUM \\
2  & \texttt{dram\_bytes\_read}            & Memory BW        & DRAM read bytes              & SUM \\
3  & \texttt{dram\_bytes\_written}         & Memory BW        & DRAM write bytes             & SUM \\
4  & \texttt{memory\_throughput\_pct}      & Memory BW        & Memory subsystem \% peak     & DW-mean \\
5  & \texttt{dram\_throughput\_pct}        & Memory BW        & DRAM \% of peak BW           & DW-mean \\
6  & \texttt{mem\_busy\_pct}              & L1 pipeline      & L1 pipeline \% peak active   & DW-mean \\
7  & \texttt{l1\_hit\_rate}               & Cache            & L1 cache hit rate (\%)       & DW-mean \\
8  & \texttt{l2\_hit\_rate}               & Cache            & L2 cache hit rate (\%)       & DW-mean \\
9  & \texttt{sm\_throughput\_pct}         & Compute          & SM \% of peak throughput     & DW-mean \\
10 & \texttt{tensor\_core\_active\_pct}   & Tensor Cores     & TC pipeline active (\%)      & DW-mean \\
11 & \texttt{sm\_active\_cycles}          & Compute          & Total SM active cycles       & SUM \\
12 & \texttt{achieved\_occupancy}         & Occupancy        & Warps active \% of peak      & DW-mean \\
13$\dagger$ & \texttt{warp\_cycles\_per\_instr}    & Latency          & Cycles per issued instruction & DW-mean \\
14 & \texttt{eligible\_\allowbreak cycles\_pct}       & Latency          & Eligible warp issue rate (\%) & DW-mean \\
15 & \texttt{executed\_instructions}      & Throughput       & Total executed instructions  & SUM \\
16 & \texttt{ipc\_active}                 & Throughput       & IPC when SM active           & DW-mean \\
17 & \texttt{avg\_threads\_per\_warp}     & Divergence       & Avg threads active per warp  & DW-mean \\
18 & \texttt{registers\_per\_thread}      & Reg.\ pressure   & Registers allocated / thread & MAX \\
19 & \texttt{local\_memory\_spills}       & Reg.\ spills     & Spill events to L1 local mem & SUM \\
20 & \texttt{dynamic\_smem\_per\_block}   & SMEM pressure    & Dynamic shared mem / block   & MAX \\
\bottomrule
\end{tabular}
\end{table}

\subsection{Prototype Automation via Claude Code Plugin}
\label{sec:appendix:plugin}

As a prototype implementation, \sys includes a Claude Code plugin that automates the
closed-loop protocol via a single \texttt{/optimize workload.py} command.  The plugin orchestrates five
specialized LLM agents---\textbf{capture-agent}, \textbf{optimization-strategist},
\textbf{backend-engineer}, \textbf{validation-agent}, and \textbf{/report}---covering all
pipeline stages from baseline capture through cross-profile comparison. 


\subsection{Annotated \texttt{profile.json} Schema}
\label{sec:appendix:schema}

\texttt{profile.json} is the central output of \sys: a JSON document that maps each
dispatched PyTorch operator to its constituent CUDA kernels, per-kernel \nsys durations,
and aggregated hardware metrics.  \tabref{tab:schema_fields} describes every top-level field; the listing below
shows a representative single-operator excerpt from the GPT-2 baseline
(\texttt{aten::mm}, 512$\times$768 $\times$ 768$\times$768 weight projection, NVTX-attributed).
The \texttt{aggregated} block condenses per-kernel counter values into a single
operator-level record using the duration-weighted mean (\equref{eq:dwmean}); these are the
values consumed by downstream analysis tools and the Claude Code plugin.
The \texttt{duration\_ns} and \texttt{aggregated.total\_duration\_ns} fields are
\nsys-derived timing fields; \ncu metrics, including
\texttt{gpu\_\_time\_duration.sum}, are stored only inside \texttt{metrics.raw}.

\begin{lstlisting}[basicstyle=\scriptsize\ttfamily, breaklines=true, frame=single,
                   xleftmargin=2pt, xrightmargin=2pt]
{
  "schema_version": "1.0",
  "capture_metadata": {
    "model_name": "gpt2",
    "device_name": "NVIDIA RTX PRO 6000 Blackwell",
    "compile_mode": "inductor",
    "capture_timestamp_utc": "2026-06-01T00:53:25Z"
  },
  "operators": [
    {
      "operator_id": "aten::mm_26",
      "operator_name": "aten::mm",
      "call_index": 26,
      "is_fused": false,
      "fused_with": [],
      "kernels": [
        {
          "kernel_id": "k_00880",
          "kernel_name": "cutlass_80_simt_sgemm_128x32_8x5_nn_align1",
          "attribution_method": "nvtx",
          "confidence": "medium",
          "duration_ns": 32736,
          "stream_id": 7,
          "grid_dim": [128, 1, 5],
          "block_dim": [128, 1, 1],
          "metrics": { "...": "20 raw ncu counter fields (architecture-specific names)" }
        }
      ],
      "aggregated": {
        "total_duration_ns": 32736,
        "kernel_count": 1,
        "tensor_core_active_pct": 0.0,
        "achieved_occupancy": 18.26,
        "sm_throughput_pct": 37.6,
        "dram_throughput_pct": 7.16,
        "l2_hit_rate": 89.55,
        "registers_per_thread": 80.0,
        "eligible_cycles_pct": 40.6,
        "local_memory_spills": 0,
        "warp_cycles_per_instruction": null
      }
    }
  ],
  "unattributed_kernels": []
}
\end{lstlisting}

Key fields: \texttt{attribution\_method} records the tier that assigned operator identity
(\texttt{"nvtx"}, \texttt{"inductor\_fusion"}, or \texttt{"torch\_profiler\_correlation"});
\texttt{confidence} mirrors the tier level (\texttt{"high"} or \texttt{"medium"});
\texttt{is\_fused} and \texttt{fused\_with} list all \texttt{aten::} ops fused into the
same Triton kernel.  The \texttt{null} value for \texttt{warp\_cycles\_per\_instruction}
reflects the absent Blackwell counter.
The \texttt{metrics.raw} sub-object stores all architecture-specific ncu column names
exactly as emitted by \texttt{ncu --import --csv}; the \texttt{aggregated} block outlines an operator-level metric summary of the attributed kernels.

\begin{table}[h]
\centering
\caption{Top-level \texttt{profile.json} field reference.}
\label{tab:schema_fields}
\scriptsize
\setlength{\tabcolsep}{4pt}
\begin{tabular}{lp{3.5cm}p{7.5cm}}
\toprule
\textbf{Field} & \textbf{Type} & \textbf{Meaning} \\
\midrule
\texttt{schema\_version}       & string        & Schema revision for forward-compatibility checks \\
\texttt{capture\_metadata}     & object        & Device name, torch/CUDA versions, nsys/ncu report paths \\
\texttt{operators[]}           & array         & One entry per attributed operator dispatch \\
\quad\texttt{operator\_id}     & string        & Unique key: \texttt{aten::op\_N} \\
\quad\texttt{operator\_name}   & string        & Canonical \texttt{aten::} name (with size annotations) \\
\quad\texttt{call\_index}      & int           & Dispatch ordinal within the measured window \\
\quad\texttt{is\_fused}        & bool          & True if kernel spans $>$1 aten op (Inductor fusion) \\
\quad\texttt{fused\_with}      & string[]      & Co-fused \texttt{aten::} op names \\
\quad\texttt{kernels[]}        & array         & Per-kernel records with raw ncu metrics \\
\quad\texttt{aggregated}       & object        & Duration-weighted operator-level metric summary \\
\texttt{unattributed\_kernels[]} & array       & Kernels matching no attribution path \\
\bottomrule
\end{tabular}
\end{table}

\subsection{Supplementary Figures}
\label{sec:appendix:candidates}

\begin{figure}[H]
\centering
\begin{tikzpicture}[font=\small, >=Stealth,
  stage/.style={draw, rounded corners=4pt, fill=white,
                minimum width=3.6cm, minimum height=0.75cm,
                align=center, inner sep=5pt, font=\small},
  hook/.style={draw, rounded corners=3pt, align=left,
               inner sep=5pt, font=\scriptsize},
  arr/.style={-Stealth,thick},
  harr/.style={-Stealth, thick, dashed},
]

  % Pipeline stages (vertical, left column)
  \node[stage, fill=gray!10] (in)     at (0,0)    {\texttt{workload.py}};
  \node[stage, fill=blue!8]  (dynamo) at (0,-2.2)  {Dynamo FX\\graph capture};
  \node[stage, fill=blue!8]  (aot)    at (0,-4.4)  {AOTAutograd\\decomposition};
  \node[stage, fill=blue!8]  (ind)    at (0,-6.6)  {Inductor\\code generation};
  \node[stage, fill=gray!10] (out)    at (0,-8.8)  {Triton / CUDA kernels};

  \draw[arr] (in)     -- (dynamo);
  \draw[arr] (dynamo) -- (aot);
  \draw[arr] (aot)    -- (ind);
  \draw[arr] (ind)    -- (out);

  % Level 1 hook (between Dynamo and AOT)
  \node[hook, fill=red!8, draw=red!50] (l1) at (5.5,-2.2)
    {\textbf{Level 1} (pre-decomposition)\\
     \textbullet\ QKV fusion (\textit{SDPA case study})\\
     \textbullet\ Attention canonicalization};
  \draw[harr, red!60] (l1.west) -- (dynamo.east);

  % Level 2 hook (between AOT and Inductor)
  \node[hook, fill=orange!8, draw=orange!50] (l2) at (5.5,-4.4)
    {\textbf{Level 2} (Aten IR)\\
     \textbullet\ BF16 dtype promotion (\textit{all case studies})\\
     \textbullet\ Weight folding, layout annotation};
  \draw[harr, orange!60] (l2.west) -- (aot.east);

  % Level 3 hook (at Inductor)
  \node[hook, fill=green!8, draw=green!50!black] (l3) at (5.5,-6.6)
    {\textbf{Level 3} (Inductor config)\\
     \textbullet\ Weight freezing (\textit{all case studies})\\
     \textbullet\ Autotuning mode selection};
  \draw[harr, green!60!black] (l3.west) -- (ind.east);

  % Compose label
  \node[font=\scriptsize, gray, align=center] at (0,-9.6)
    {Passes at each level compose\\without interfering across levels};

\end{tikzpicture}
\caption{\textbf{Three-level FX pass hook points.}
  The \texttt{torch.compile()} pipeline lowers a model through three IR stages.
  Level~1 passes must run before AOTAutograd because structural patterns such as
  a shared \texttt{LayerNorm} output feeding Q, K, and V projections---the
  precondition for QKV fusion---are decomposed into per-consumer copies by
  AOTAutograd and become invisible afterward.  Level~2 (Aten IR) applies
  dtype and layout transforms on the fully decomposed graph.  Level~3
  Inductor flags control non-graph behaviors such as weight freezing and
  autotuning that cannot be expressed as graph rewrites.}
\label{fig:cand:fxpasslevel}
\end{figure}


\end{document}
